% Copyright (c) 2026 Geovana Neves. Licensed under CC BY 4.0 (https://creativecommons.org/licenses/by/4.0/). See LICENSE-AND-AUTHORSHIP.md.
%
\section{The unsteady POLAR}

\begin{frame}{The unsteady POLAR averages the plots history}
\begin{itemize}\small
  \item The source is the plots history, time-averaged over the row's window.
        It does not read the native coefficient export for that average.
  \item Columns are plot variables under the names the export prints,
        rather than the steady polar's fixed twenty-four columns.
  \item Flight-condition and reference-length columns are added from the workspace.
  \item The native export states only the last time step, one instant of a
        rotor cycle. It still ships as a health check.
\end{itemize}
\end{frame}

\begin{frame}{One file per simulation, one row per point}
\small \code{polars/P<sim>\_<name>\_uns\_avg.csv}
\begin{itemize}
  \item Every file under \code{post/} comes from a sweep.
        The name states that this is an average of unsteady history.
  \item It carries the \code{P} used by every per-point product.
\end{itemize}
\end{frame}

\begin{frame}{Each polar row states its window and moment point}
\begin{enumerate}
  \item \code{FIRST\_STEP}, \code{LAST\_STEP}, \code{STEPS} open the row
        and identify that point's averaging window.
  \item The condition block follows.
  \item Then \code{XMOM}, \code{YMOM}, \code{ZMOM} state the moment point.
\end{enumerate}
\begin{takeaway}
The plots carry moments. A moment needs its point.
\end{takeaway}
\end{frame}

\begin{frame}{The super-file content joins the polar}
\small After the plot columns, the table adds:
\begin{itemize}
  \item The matrix row's cells and the run record's scalars.
  \item Each rotor's speed.
  \item The solver flags.
\end{itemize}
\begin{takeaway}
For an unsteady point, this content is part of the polar, not a second super file.
\end{takeaway}
\end{frame}

\begin{frame}{Each global plot group supplies eighteen axis columns}
\begin{itemize}\small
  \item The axis coefficients follow the plot columns, using the steady
        polar's names in this order:
        \code{CDW .. CNW25}, \code{CDS .. CNS25}, \code{CDB .. CNB25}.
  \item Each name takes the plot group's \textbf{whole} name as a suffix,
        for example \code{CLW\_MRP\_TOTAL}.
  \item A second global-frame group adds its own eighteen columns
        and renames none.
\end{itemize}
\end{frame}

\begin{frame}{The unsteady axes use dimensional global loads}
\begin{itemize}\small
  \item Use \code{FX, FY, FZ, MX, MY, MZ} from a pproc plot group with
        \code{frame = "MRP"}, in Newtons and Newton metres.
  \item Average these components over the row's window.
        Divide by that row's
        $\tfrac12\mathrm{RHO}\,\mathrm{VINF}^2\,\mathrm{SREF}$
        and also by \code{CREF} for moments.
  \item Turn them as the steady-polar axes are turned, with the same
        force signs and final moment normalisation.
  \item A rotor's own frame cannot supply this block: its axes are not
        the geometry's axes.
\end{itemize}
\end{frame}

\begin{frame}{Missing global plots and the axes block}
\begin{itemize}\small
  \item If the pproc has no global-frame group plotting all six components,
        the run adds \code{MRP\_TOTAL} over every boundary, where the run
        has an \code{MRP} frame.
  \item A pproc that already plots them stays unchanged, to the byte.
  \item An older run without those plots has no axes block.
        \code{products.json} records this under \code{polars/<file>\#axes}.
\end{itemize}
\begin{takeaway}
Missing source plots do not produce an axes block filled with \code{NA}.
\end{takeaway}
\end{frame}

\begin{frame}[fragile]{The pproc can declare output column names}
\small The package cannot infer which plot label carries which coefficient.
Undeclared names pass through exactly as the export prints them.
\begin{lstlisting}[style=ftsbase]
[names]
CL_MRP_TOTAL = "CL_TOTAL"
FX_HUB_PUSHER = "FX_PUSHER"
\end{lstlisting}
\small This dictionary maps an exported plot column to the name
a downstream reader wants. Inventing a source label would otherwise risk
an entire column of \code{NA}.
\end{frame}

\begin{frame}{The names dictionary changes selected headings last}
\begin{itemize}\small
  \item It renames the unsteady polar and averaged reductions:
        \code{time\_average} and the passage series.
  \item \code{probes/<point>\_plots.csv} keeps the export's names because
        it is the source table.
  \item Per-blade and azimuthal tables are not renamed: their columns are
        no longer the export's own names.
  \item Axes and \code{[equations]} read export names.
        The dictionary applies last, to the heading alone.
\end{itemize}
\end{frame}

\begin{frame}{The whole dictionary applies, or none of it does}
\begin{itemize}\small
  \item A source column that no plot of the point prints invalidates the
        dictionary for that point.
  \item A destination name already carried by the table does the same.
  \item Every column then retains the export's name, with a warning and
        \code{products.json} entries under \code{polars/<file>\#names}
        and \code{probes/<point>\#names}.
\end{itemize}
\begin{takeaway}
A failed rename never becomes a column of \code{NA}.
\end{takeaway}
\end{frame}

\begin{frame}{Names that are occupied or refused}
\begin{itemize}\small
  \item \code{CL} is taken. The native export's last-step \code{CL},
        \code{CDi}, \code{CDo}, \code{Cx} and the rest already appear
        in the polar's setup content as a health check.
  \item A plotted column cannot take one of those names.
        \code{CL\_TOTAL} is available instead.
  \item Two dictionary entries giving one name, or a name that is not
        one word, are refused when the pproc is read.
\end{itemize}
\end{frame}
